\documentclass[a4j,dvipdfmx]{jsarticle}
\usepackage{amsmath,amssymb}
\usepackage{siunitx}
\usepackage{ascmac}
\usepackage{amsthm}
\usepackage{bm}

\usepackage[margin=10truemm]{geometry}

\newcommand{\R}{\mathbb{R}}
\newcommand{\rank}{\mathrm{rank}}
\newcommand{\tr}{\mathrm{tr}}
\newcommand{\Ker}{\mathrm{Ker}\hspace{0.5mm}}
\newcommand{\E}{\mathcal{E}}
\newcommand{\F}{\mathcal{F}}
\renewcommand{\labelenumi}{(\arabic{enumi})}
\renewcommand{\Im}{\mathrm{Im}\hspace{0.5mm}}
% \usepackage[dvipdfmx]{hyperref}
\begin{document}
    \part*{第三回 線型代数 問題}
    \subsection*{問題1}
        $(m,n)$型行列$A$の階数ははじめ「基本変形で標準形に変形した時, 対角線上に並ぶ1の数」と定義されたが, これと同値な4つの条件を答えよ.
    \subsection*{問題2}
        写像$f:M_n(\R)\to\R, A\mapsto\tr A$と定める. このとき, 次の問いに答えよ.
        \begin{enumerate}
            \item $f$が線形写像であることを示せ.
            \item $\Im f$を求めよ.
            \item $\Ker f$の次元を求めよ. 
        \end{enumerate}
    \subsection*{問題3}
        2次以下の実係数多項式空間$P_2(\R)$の線形変換$T,S$を
        \begin{align*}
            (Tf)(x)&=\int_{-1}^{1}(x-t)^2f(t)dt\\
            (Sf)(x)&=e^{x}\frac{d}{dx}(e^{-x}f(x))
        \end{align*}
        によって定義することができる. 基底$<1,x,x^2>$に関する$T,S$の表現行列を求めよ.
    \subsection*{問題4}
        $A=\begin{pmatrix}-1 & -5 & -1\\2 & 4 & 0\\-4 & -8 & 0\end{pmatrix}$とするとき, $\R^3$の元$\bm{x}$に$A\bm{x}$を対応させる写像を$T_A$とする.
        \begin{enumerate}
            \item $T_A$が$\R^3$の線形変換であることを示せ.
            \item $\F = <\begin{pmatrix}1\\-1\\2\end{pmatrix},\begin{pmatrix}3\\-2\\4\end{pmatrix},\begin{pmatrix}2\\-1\\3\end{pmatrix}>$が$\R^3$の基底であることを示せ.
            \item $\E_0 = <\begin{pmatrix}1\\0\\0\end{pmatrix},\begin{pmatrix}0\\1\\0\end{pmatrix},\begin{pmatrix}0\\0\\1\end{pmatrix}>$とする. このとき, $\E_0\to\F$の基底の取替行列$P$を求めよ.
            \item $\F$に関する$T_A$の表現行列を求めよ.
        \end{enumerate}
    \subsection*{問題5}
        数列空間$V=\{(x_n)\mid x_{n+2}-3x_{n+1}+2x_{n}=0\}$について, 以下の問に答えよ. ただし, 和を$(a_n)+(b_n)=(a_n+b_n)$, スカラー倍を$\lambda (a_n)=(\lambda a_n)$と定める.
        \begin{enumerate}
            \item 項をずらす写像$f:V\to V,(x_{n})\mapsto (x_{n+1})$は線形写像であることを示せ.
            \item $x_0=1,x_1=0$である数列を$\bm{e}_0$とし, $x_0=0,x_1=1$である数列を$\bm{e}_1$とする. この時, $\E=<\bm{e}_0,\bm{e}_1>$が基底であることを示せ.
            \item $x_0=1,x_1=1$である数列を$\bm{f}_0$とし, $x_0=1,x_1=2$である数列を$\bm{f}_1$とする. この時, $\F=<\bm{f}_0,\bm{f}_1>$が基底であることを示せ.
            \item $\E\to\F$への基底の取替え行列$P$を求めよ.
            \item $f$の$\F$に関する表現行列を求めよ.
            \item 漸化式$x_{n+2}-3x_{n+1}+2x_{n}=0$の一般解を求めよ.
        \end{enumerate}
    \clearpage
    \subsection*{問題6}
        Cauchy-Schwarzの不等式を証明せよ. また三角不等式を証明せよ.
    \subsection*{問題7}
        $\R^3$に標準内積を入れたとき, 3つのベクトル$\E = < \begin{pmatrix}1 \\ -1 \\ 0\end{pmatrix},\begin{pmatrix}1 \\ 0 \\ -1\end{pmatrix},\begin{pmatrix}1 \\ 2 \\ 3\end{pmatrix}>$は基底である. この基底からSchmidtの直交化法より正規直交基底を作れ.
    \subsection*{問題8}
        $V$を$1,\cos x,\sin x$によって張られる空間とする. これは3次元の実線型空間で, 基底はもちろん$\E=<1,\sin x,\cos x>$である. $V$内の二つの関数$f,g$に対し, その内積$(f,g)$を
        \begin{equation*}
            (f,g)=\int_{-\pi}^{\pi}f(x)g(x)dx
        \end{equation*}
        として定義する. 線形変換$T$を$T:f(x)\mapsto f(x+c)$とするとき, 以下の問いに答えよ.
        \begin{enumerate}
            \item 上の$T$が本当に$V$の線形変換であることを確かめよ. つまり, $V$の任意の元を$T$で飛ばしたときに, それがまた$V$に入っていることを確かめよ.
            \item $T$が$V$の直交変換であることを示せ.
            \item $V$の正規直交基底を一つ求めよ.
            \item (3)で求めた正規直交基底に関する$T$の表現行列を求めよ.
        \end{enumerate}
    \subsection*{問題9}
        計量空間$V$の部分空間$W$に対し, $W$のすべての元と直交するような$V$の元全体を$W$の\textbf{直交補空間}といい$W^{\bot}$で表す. このとき, 以下の問いに答えよ.
        \begin{enumerate}
            \item $W^{\bot}$は$V$の部分空間であることを示せ.
            \item $V=W\oplus W^{\bot}$, すなわち, $V$は$W$と$W^{\bot}$の\textbf{直和}であることを示せ.\footnote{一般に, ベクトル空間$V$の部分空間$W_1,W_2$について, $V$が$W_1,W_2$の和空間$W_1+W_2$で, $V$の任意のベクトルの$W_1,W_2$のベクトルの和として表す仕方が一意的であるとき, 直和であるという.}
        \end{enumerate}

    \subsection*{問題10}
        計量同型写像の合成写像は計量同型写像であることを示せ. また, 計量同型写像の逆写像も計量同型写像であることを示せ.
    
    \subsection*{補充問題}以下余裕があれば証明せよ.
        \begin{enumerate}
            \item 線形写像$f$が単射であることと$\Ker f=\{\bm{0}\}$が同値であることを示せ.
            \item 任意の線形写像$f$について, $f(v_1),\dots,f(v_k)$が一次独立なら$v_1,\dots,v_k$は一次独立であることを示せ.
            \item 単射な線形写像$f$については(2)の逆が成り立つことを示せ.
        \end{enumerate}
\end{document}